\section{Artifact and Reproducibility}

The artifact is a reusable benchmark rather than a single script: a core package (deletion targets, artifacts, probes, responses, evidence vectors, provenance graphs, clean-counterfactual distinguishability tests, risk scoring, audit algorithms, RAG/LoRA-SFT/DenseRAG harnesses, and reporting) plus scripts that generate the JSON/CSV rows, aggregate tables, figures, and manuscript PDF via a one-command run. The benchmark card specifies tracks, metrics, baselines, clean-reference construction, adversary classes, and release assumptions, and every exported row records the track, deletion/audit method, decision, failure and detection labels, risk score, leakage channels, retain utility, audit calls, seed, target size, and budget (advanced rows add pretrained-LoRA forget margins, retain perplexity/accuracy, embedding-model names, and top-$k$ deleted-derivative hit rates). A requirements file and model-cache manifest document expected runtime and cache assumptions.

To preserve double-anonymous review, the review manuscript omits identifying repository links; the artifact package itself includes the harness, experiment scripts, generated tables, manuscript sources, environment requirements, benchmark card, and redacted adjudication logs, with a de-anonymized public link to follow in the final version. The current artifact uses private data and public FEVER evidence, which avoids releasing sensitive records while still exercising realistic failure modes (index-only deletion, shadow-copy and cache retention, adapter memory, negative/gradient-ascent unlearning, -derivative retention, and over-unlearning). It should be read as a technical audit benchmark, not a legal-compliance certification tool.
